% ============================================================
% VERSIÓN INDIVIDUAL — DANIEL DUARTE CORDERO
% Baseline análogo: AAIA'17 Hearthstone.
% ============================================================
\section{Baseline de Literatura}
\label{sec:baseline}

\subsection{Origen del baseline}
\label{subsec:origenbaseline}
El dataset propio carece de un resultado arbitrado previo directamente
comparable. Se adopta como referencia análoga el AAIA'17 Data Mining
Challenge sobre Hearthstone, cuya tarea consistía en estimar la
probabilidad de victoria a partir de estados intra-partida y cuya
métrica de competencia era AUC \cite{janusz_hearthstone_2017}. El
artículo resumen reporta un baseline oficial de 0.7846 AUC y un mejor
resultado final de 0.8019 AUC. La solución ganadora utilizó un ensamble
de redes neuronales, pero ese detalle no convierte el valor en un
benchmark equivalente para Pokémon TCG.

\subsection{Número objetivo}
\label{subsec:objetivobaseline}
\begin{quote}
\textbf{Objetivo declarado:} rango de referencia AUC 0.7846--0.8019 en
el análogo de Hearthstone \cite{janusz_hearthstone_2017}. No se declara
como umbral obligatorio a superar, sino como orden de magnitud externo
para una tarea de predicción de victoria desde estado.
\end{quote}

\subsection{Comparabilidad del protocolo}
\label{subsec:comparabilidad}
La comparación es \textbf{parcialmente comparable}. Coinciden la
naturaleza binaria del desenlace, el uso de información intra-partida y
la evaluación mediante AUC, pero no coinciden el juego, las reglas, la
representación del estado ni la población que genera las partidas. El
valor de Hearthstone sirve para contextualizar el desempeño, no para
afirmar superioridad directa.

\noindent\textbf{Diferencias que impiden una comparación justa.}
\begin{itemize}
  \item La partición y el mecanismo de generación de ejemplos del reto
  de Hearthstone no son equivalentes a nuestro split agrupado por
  \texttt{episode\_id} ni al protocolo temporal.
  \item Aunque ambos reportan AUC, nuestro estudio añade log loss,
  Brier score y calibración, por lo que la evaluación no se reduce a
  la misma métrica.
  \item Los estados, cartas, recursos, reglas y representación de
  Hearthstone son distintos de los snapshots observables de Pokémon
  TCG.
  \item La población de bots, mazos y distribución competitiva difiere
  de los agentes artificiales y mazos de la competencia Pokémon TCG.
\end{itemize}

\noindent\textbf{Qué significa entonces el número.} El intervalo
0.7846--0.8019 se usa como referencia externa del nivel de
discriminación alcanzado en una tarea análoga. No es una cifra que el
proyecto deba ``batir'' para considerarse válido: un resultado inferior
puede ser metodológicamente sólido si se obtiene bajo un protocolo más
estricto, y uno superior no probaría automáticamente mejor desempeño
por las diferencias anteriores.
